\begin{ex}
Số liền trước của 800 là số nào?
\shortans{799}
\loigiai{
$800 - 1 = 799$.
}
\end{ex}

\begin{ex}
Số liền sau của 699 là số nào?
\shortans{700}
\loigiai{
$699 + 1 = 700$.
}
\end{ex}

\begin{ex}
Tính: $356 + 243 = \dots$
\shortans{599}
\loigiai{
$356 + 243 = 599$.
}
\end{ex}

\begin{ex}
Tính: $785 - 342 = \dots$
\shortans{443}
\loigiai{
$785 - 342 = 443$.
}
\end{ex}

\begin{ex}
Tính: $124 \times 3 = \dots$
\shortans{372}
\loigiai{
$124 \times 3 = 372$.
}
\end{ex}

\begin{ex}
Tính: $648 : 2 = \dots$
\shortans{324}
\loigiai{
$648 : 2 = 324$.
}
\end{ex}

\begin{ex}
Tính giá trị biểu thức: $40 + 60 : 2$.
\shortans{70}
\loigiai{
$40 + 30 = 70$.
}
\end{ex}

\begin{ex}
Tính giá trị biểu thức: $(50 + 30) \times 3$.
\shortans{240}
\loigiai{
$80 \times 3 = 240$.
}
\end{ex}

\begin{ex}
Số 42 gấp số 7 mấy lần?
\shortans{6}
\loigiai{
$42 : 7 = 6$ lần.
}
\end{ex}

\begin{ex}
Số 72 gấp số 9 mấy lần?
\shortans{8}
\loigiai{
$72 : 9 = 8$ lần.
}
\end{ex}

\begin{ex}
Tìm $\frac{1}{6}$ của 54 quả cam.
\shortans{9}
\loigiai{
$54 : 6 = 9$ quả cam.
}
\end{ex}

\begin{ex}
Tìm $\frac{1}{8}$ của 48 kg gạo.
\shortans{6}
\loigiai{
$48 : 8 = 6$ kg gạo.
}
\end{ex}

\begin{ex}
Một hình chữ nhật có chiều dài 15 cm, chiều rộng 5 cm. Chu vi hình chữ nhật đó là bao nhiêu xăng-ti-mét?
\shortans{40}
\loigiai{
$(15 + 5) \times 2 = 40$ cm.
}
\end{ex}

\begin{ex}
Một hình vuông có cạnh 9 cm. Chu vi hình vuông đó là bao nhiêu xăng-ti-mét?
\shortans{36}
\loigiai{
$9 \times 4 = 36$ cm.
}
\end{ex}

\begin{ex}
Hình tròn có bán kính 7 cm thì đường kính là bao nhiêu xăng-ti-mét?
\shortans{14}
\loigiai{
$7 \times 2 = 14$ cm.
}
\end{ex}

\begin{ex}
Một thùng chứa 120 lít dầu, người ta bán đi $\frac{1}{3}$ số dầu đó. Hỏi đã bán bao nhiêu lít dầu?
\shortans{40}
\loigiai{
$120 : 3 = 40$ lít.
}
\end{ex}

\begin{ex}
Một cửa hàng ngày thứ nhất bán được 135 kg gạo, ngày thứ hai bán gấp đôi ngày thứ nhất. Cả hai ngày bán được bao nhiêu ki-lô-gam gạo?
\shortans{405}
\loigiai{
Ngày 2: $135 \times 2 = 270$ kg. Cả hai ngày: $135 + 270 = 405$ kg.
}
\end{ex}

\begin{ex}
Khối lập phương có bao nhiêu đỉnh?
\shortans{8}
\loigiai{
Có 8 đỉnh.
}
\end{ex}

\begin{ex}
Khối hộp chữ nhật có bao nhiêu mặt?
\shortans{6}
\loigiai{
Có 6 mặt.
}
\end{ex}

\begin{ex}
Đổi: $5\text{ cm } 6\text{ mm} = \dots\text{ mm}$.
\shortans{56}
\loigiai{
$50 + 6 = 56$ mm.
}
\end{ex}

\begin{ex}
Đổi: $1\text{ kg} - 250\text{ g} = \dots\text{ g}$.
\shortans{750}
\loigiai{
$1000 - 250 = 750$ g.
}
\end{ex}

\begin{ex}
Đổi: $1\text{ l} - 450\text{ ml} = \dots\text{ ml}$.
\shortans{550}
\loigiai{
$1000 - 450 = 550$ ml.
}
\end{ex}

\begin{ex}
Nhiệt độ trong phòng học là $25^\circ$C, ngoài trời là $33^\circ$C. Ngoài trời nóng hơn trong phòng bao nhiêu độ C?
\shortans{8}
\loigiai{
$33 - 25 = 8^\circ$C.
}
\end{ex}

\begin{ex}
Tìm $x$: $x \times 4 = 360$.
\shortans{90}
\loigiai{
$x = 360 : 4 = 90$.
}
\end{ex}

\begin{ex}
Tìm $x$: $x : 3 = 120 + 10$.
\shortans{390}
\loigiai{
$x : 3 = 130 \Rightarrow x = 130 \times 3 = 390$.
}
\end{ex}
